\begin{proof}
\begin{enumerate}[label=\arabic*.]
\item Use the explicit matrix $C_2$ in the proof of
\cref{maps:src:F.2.6}.  In the original vertical marking $D$ is the
bottom-right block of $M_2$; multiplication gives
$DC_2=C_2D_2$, where $D_2$ is the four-cycle followed by the sign
on its fifth coordinate.  The change of cochain coordinates is
$\bigwedge^kC_2^t$, and the inverse is $\bigwedge^kC_2^{-t}$.
Thus the product cell structure of the signed-permutation torus,
transported by this integral torus isomorphism, has zero cellular
differential and the stated action $T=\bigwedge^{q-1}D^t$.
The fibre cells and their prisms give the actual mapping-torus
cochain complex.  The degree-$q$ suspension classes are consequently
the coinvariants of $T$ on vertical degree-$q-1$ cochains.

\item Write $u=-x_0$ in the original marking.  The affine action is
$(u,y,z)\mapsto(u+1/4,Dy,iz)$, and
$W=e^{-2\pi i u}z$ is invariant.  With core angle $a=4u\bmod1$
and normal angle $k=\arg W/(2\pi)$, the coarse angle is $a+4k$.
Let $g,h$ denote the two corresponding loops.  The integral basis
\[
 v=g^4h^{-1},\qquad t=g,
 \qquad [v,t]_{[g,h]}=\begin{pmatrix}4&1\\-1&0\end{pmatrix}
\]
has determinant one.  Its coarse degrees are zero and one; its
vertical monodromies are $T^4=I$ and $T$ respectively.
It specifies the actual generic six-torus fibre and the primitive
base loop in the boundary product.

\item A suspension cochain representing $[w]$ has values $(w,0)$
on $(g,h)$.  Following the four consecutive core prisms transports
their values by $I,T,T^2,T^3$, giving
\[
 (w,0)_{(g,h)}\longmapsto(N_Tw,w)_{(v,t)},
 \qquad N_T=I+T+T^2+T^3.
 \tag{FT.1}\label{maps:eq:FT.1}
\]
For a torsion coinvariant there are an integer $n\ne0$ and an
integral cochain $b$ with $nw=(T-I)b$.  Therefore
$nN_Tw=N_T(T-I)b=0$, and $N_Tw=0$ in the free cochain lattice.
The fibre restriction vanishes; its base-circle component is $[w]$.
Replacing $w$ by $w+(T-I)b$ changes that component by its cellular
coboundary.  The formula is consequently a well-defined map of the
stated integral coinvariant groups.

\item The base cover by a disc at the finite point and its
complement gives the relative \v{C}ech connecting map.  On their
annular intersection the degree-zero fibre restriction is zero and
the next component is the one-cocycle just computed.
The connecting map sends it to $\pm[w]$ in the global coinvariants;
the two ordinary peripheral loops impose the relations
$\bigwedge^{q-1}M_1^t-I$ and $\bigwedge^{q-1}M_2^t-I$.
These are the two successive boundaries in the same base--fibre
complex, so their composition is the Leray $d_2$ of
\eqref{maps:eq:F.6}.  The one boundary-orientation choice accounts
for the sign.  Its coefficient is one, as witnessed by the
degree-one loop $t$ above.  Point-supported kernels or cokernels in
the comparison of restriction sheaves have no positive cohomology,
so they preserve this degree-two base quotient.

\item In signed-permutation coordinates the negative orbits in
vertical degrees $1,2,3$ are represented respectively by
$e_4$, $e_0\wedge e_2$ and $e_0\wedge e_2\wedge e_4$.
The orbit relations give exactly one local $\mathbb Z/2$ in each
case.  Transport by $\bigwedge^kC_2^{-t}$ and then include vertically
in the six-lattice.  This is the same map as the vertical inclusion
followed by $\bigwedge^kQ_2$ used in
\cref{maps:src:F.2.6}, because $P_2=(-1)\oplus C_2$.
The explicit $w_k,\rho_k,E_k,z_k$ there satisfy
\[
 \rho_k\Delta_k=0\pmod2,\qquad
 \rho_kw_k=1\pmod2,\qquad E_kz_k=2w_k,
 \qquad |\det E_k|=2.
\]
Thus the unique nonzero local class maps to the unique nonzero
global class in each of the three degrees.  Formula
\eqref{maps:eq:FT.1} sends the zero representative to zero as well.
\end{enumerate}
\end{proof}
